% =====================================================================================
% SMS v2 — Memoria persistente en agentes LLM para ingeniería de software
% Protocolo: research/SMS-V2-PROTOCOL.md · Ley de cierre §9 · Ledger de números: claims.yml
%
% REGLA DURA (§9.5): ninguna cifra se escribe a mano. Todo número de resultados es una macro
% \cl... definida en claims.tex, que verify_sms_claims.py re-deriva del corpus committeado.
% Dentro de las regiones %%CLAIMS-STRICT el gate RECHAZA cualquier dígito literal.
% =====================================================================================
\documentclass[conference]{IEEEtran}

\usepackage[utf8]{inputenc}
\usepackage[T1]{fontenc}
\usepackage[spanish,es-tabla,es-noquoting]{babel}
\usepackage{booktabs}
\usepackage{tabularx}
\usepackage{array}
\usepackage{amssymb}
\usepackage{graphicx}
\usepackage{xcolor}
\usepackage{multirow}
\usepackage[hidelinks]{hyperref}
\usepackage{url}

% GENERATED by harvest/verify_sms_claims.py — DO NOT EDIT BY HAND.
% Every macro below is re-derived from the committed corpus (protocol §9.5).
% If you want to change a number, change the corpus or the derivation, never this file.

% clCorpusUnicos: ¿Cuántos registros únicos dejó la búsqueda académica tras dedup?
\newcommand{\clCorpusUnicos}{1.069}
% clIncluidosPrincipal: ¿Cuántos incluyó la pasada principal de cribado?
\newcommand{\clIncluidosPrincipal}{422}
   % <- GENERADO. Todos los números del paper viven aquí.

\newcolumntype{R}{>{\raggedleft\arraybackslash}X}
\newcommand{\facet}[1]{\textsc{\small #1}}

\begin{document}

\title{Memoria persistente en agentes LLM para ingeniería de software:\\
un mapeo sistemático de \clCarAnalizados{} estudios con verificación automatizada}

\author{\IEEEauthorblockN{Catalina Lazo Locky}
\IEEEauthorblockA{Departamento de Ciencias de la Computación e Informática\\
Universidad de La Frontera, Temuco, Chile\\
catalinalazolocky@gmail.com}}

\maketitle

\begin{abstract}
% ABSTRACT — se escribe AL FINAL, cuando los números están cerrados.
\end{abstract}

\begin{IEEEkeywords}
mapeo sistemático, agentes LLM, memoria persistente, ingeniería de software, gestión de contexto,
recuperación de información, evaluación de agentes
\end{IEEEkeywords}

% =====================================================================================
\section{Introducción}
\label{sec:intro}

Los agentes basados en modelos de lenguaje dejaron de ser demostraciones de una sola respuesta para
convertirse en procesos que operan durante días sobre un repositorio real. En ese régimen la
restricción dominante ya no es generar código correcto en un intento, sino \emph{sostener la
continuidad}: cada sesión nueva empieza sin recuerdo de la anterior, y lo que el agente sabe de un
proyecto es exactamente lo que alguien decidió escribir en alguna parte antes de que la ventana de
contexto se cerrara. El problema de ingeniería que sigue abierto no es cuánto sabe el modelo, sino
qué se persiste, cuándo se escribe, qué se vuelve a leer y cómo se sabe que lo persistido todavía es
verdad.

La literatura responde a esa pregunta con una palabra sobrecargada. «Memoria» nombra a la vez el
borrador intra-tarea que se descarta al terminar, el resumen que sobrevive a una compactación, el
índice vectorial de conversaciones pasadas, el grafo de entidades de un proyecto y los pesos
actualizados en tiempo de prueba. Bajo el mismo término conviven mecanismos con propiedades
incompatibles, y las decisiones de diseño —¿índice semántico o lectura directa de archivos?,
¿escritura discrecional del modelo o disparada por un evento?— se toman sin un mapa de qué se ha
medido y con qué evidencia.

%%CLAIMS-STRICT
Un mapeo sistemático previo del equipo abordó esta área con un corpus de veintisiete estudios
anclado en depuración y reparación automática de programas. Ese anclaje quedó estrecho: el objeto de
estudio se desplazó desde la tarea puntual hacia la operación continua sobre un proyecto. Este
trabajo rehace el mapeo desde cero con un ancla distinta —memoria de proyecto en agentes de larga
vida— y a otra escala: \clSelUnicos{} registros únicos cribados, \clSelIncluidosTotal{} estudios
incluidos y \clCarAnalizados{} analizados contra su texto completo, de los cuales \clCarRecientesPct{}\,\%
se publicaron en los últimos dos años.
%%CLAIMS-STRICT-END

El estudio hace dos aportes. El primero es el mapa: siete preguntas de investigación que cubren
tipos y representación, ciclo de vida de escritura, recuperación e instrumentación, persistencia
cross-sesión y multi-agente, evaluación y costo, riesgos, y práctica industrial de los proveedores de
modelos. El segundo es metodológico. Un corpus de este tamaño no se lee a mano en el marco de un
trabajo de titulación, y delegarlo a un modelo de lenguaje sin más introduce exactamente los riesgos
que la comunidad ya documentó: números que no existen, citas que no resuelven y clasificaciones que
nadie puede auditar. La respuesta que aquí se propone no es confiar en el modelo ni renunciar a él,
sino \emph{acotar cada decisión delegada con una compuerta determinista}: vocabularios cerrados que
un validador rechaza, aritmética de rúbrica recomputada por máquina, una segunda lectura ciega que
produce un coeficiente de acuerdo, y un registro de afirmaciones donde cada cifra del artículo se
re-deriva del corpus versionado en cada compilación. Ninguna de esas compuertas juzga la calidad de
la prosa; todas hacen imposible una clase concreta de error.

Finalmente, el mapeo está \emph{instrumentado para una decisión}. Las preguntas no se eligieron por
completitud temática sino porque cada una corresponde a una decisión de diseño ya tomada y
congelada en un sistema real —el prototipo Meta-Factory, un agente de larga vida cuya memoria son
documentos versionados—, y cada estudio analizado se clasificó según si la confirma, la extiende, es
neutral o la refuta. Esa última categoría es la más valiosa y la que organiza la discusión: obliga a
buscar activamente la evidencia que contradice las decisiones propias en vez de la que las
acompaña.

% =====================================================================================
\section{Antecedentes}
\label{sec:background}

\subsection{Agentes LLM de larga vida en ingeniería de software}
\subsection{Memoria, contexto y la distinción que este estudio sostiene}
\subsection{Tipos de memoria heredados de la ciencia cognitiva}

% =====================================================================================
\section{Trabajos relacionados}
\label{sec:related}

% =====================================================================================
\section{Metodología}
\label{sec:method}

\subsection{Diseño del estudio y preguntas de investigación}
\label{sec:rqs}

\subsection{Estrategia de búsqueda y fuentes}
\label{sec:search}

\subsection{Criterios de selección y rúbrica}
\label{sec:criteria}

\subsection{Proceso de selección}
\label{sec:selection}

\subsection{Snowballing y conjunto de control}
\label{sec:snowball}

\subsection{Extracción y síntesis tipada}
\label{sec:extraction}

\subsection{Fiabilidad: doble lectura independiente}
\label{sec:reliability}

\subsection{Automatización verificada: qué decidió la máquina, qué decidió un modelo, qué decidió una persona}
\label{sec:automation}

% =====================================================================================
\section{Resultados}
\label{sec:results}

%%CLAIMS-STRICT
Todos los resultados se computan sobre los \clCarAnalizados{} estudios analizados a texto completo.
Los campos son multi-etiqueta: un estudio puede declarar varios tipos de memoria o varios mecanismos
de recuperación, de modo que las columnas suman más que el total. Las proporciones se calculan
siempre sobre \clCarAnalizados{}.

\subsection{Caracterización del corpus}
\label{sec:res-corpus}

El corpus es notablemente joven: \clCarRecientes{} de los \clCarAnalizados{} estudios
(\clCarRecientesPct{}\,\%) se publicaron en los dos últimos años. Del total incluido
(\clSelIncluidosTotal{}), se obtuvo texto completo de \clSelConTexto{} (\clSelCobertura{}\,\%). El
escaneo de integridad detectó \clCalDuplicados{} pares casi-duplicados entre los incluidos, y
\clCalArtefactos{} registros resultaron ser herramientas o artefactos de software en lugar de
estudios, catalogados aparte.

\subsection{RQ1 · Tipos de memoria y sustratos de representación}
\label{sec:rq1}

La memoria episódica y la semántica dominan, y la combinación es la norma: \clRQunoHibridosTipos{}
estudios (\clRQunoHibridosPct{}\,\%) declaran dos o más tipos, frente a \clRQunoTiposUno{} que
declaran uno solo. El polo paramétrico —escribir la experiencia en los pesos del modelo— está
prácticamente desierto, y \clRQunoSinMecanismo{} estudios no aportan mecanismo propio porque su
contribución es un benchmark, un dataset o una medición.

\begin{table}[!t]
\caption{RQ1 · Tipos de memoria y sustrato de representación (multi-etiqueta, $n=\clCarAnalizados$)}
\label{tab:rq1}
\centering
\begin{tabularx}{\columnwidth}{Xr@{\hskip 1.2em}Xr}
\toprule
\textbf{Tipo de memoria} & \textbf{n} & \textbf{Representación} & \textbf{n} \\
\midrule
Episódica      & \clRQunoEpisodica     & Híbrida          & \clRQunoRepHibrida \\
Semántica      & \clRQunoSemantica     & Documentos/archivos & \clRQunoRepDocs \\
Procedimental  & \clRQunoProcedimental & Vectorial        & \clRQunoRepVector \\
De trabajo     & \clRQunoTrabajo       & Grafo            & \clRQunoRepGrafo \\
Multimodal     & \clRQunoMultimodal    & Parámetros       & \clRQunoRepParametros \\
Paramétrica    & \clRQunoParametrica   & (no declara)     & \clRQunoRepNa \\
(sin mecanismo)& \clRQunoSinMecanismo  &                  & \\
\bottomrule
\end{tabularx}
\end{table}

\subsection{RQ2 · Ciclo de vida de escritura y qué lo dispara}
\label{sec:rq2}

La escritura es la operación universalmente presente (\clRQdosWrite{}), pero el ciclo de vida
completo no lo es: la actualización de una memoria existente aparece en \clRQdosActualizacion{}, la
consolidación en \clRQdosConsolidacion{}, la admisión —decidir qué merece entrar— en
\clRQdosAdmision{} y el olvido en \clRQdosOlvido{}. El hallazgo con mayor consecuencia de diseño está
en el disparador: entre los estudios que declaran uno, \clRQdosDeterministaPct{}\,\% usan un
mecanismo no discrecional.

\begin{table}[!t]
\caption{RQ2 · Operaciones de ciclo de vida y disparador de escritura ($n=\clCarAnalizados$)}
\label{tab:rq2}
\centering
\begin{tabularx}{\columnwidth}{Xr@{\hskip 1.2em}Xr}
\toprule
\textbf{Operación} & \textbf{n} & \textbf{Disparador} & \textbf{n} \\
\midrule
Escritura      & \clRQdosWrite          & Evento determinista & \clRQdosTriggerDeterminista \\
Actualización  & \clRQdosActualizacion  & Discrecional del modelo & \clRQdosTriggerDiscrecional \\
Consolidación  & \clRQdosConsolidacion  & Presupuesto        & \clRQdosTriggerPresupuesto \\
Admisión       & \clRQdosAdmision       & Checkpoint         & \clRQdosTriggerCheckpoint \\
Olvido         & \clRQdosOlvido         & (no declara)       & \clRQdosTriggerNa \\
(no declara)   & \clRQdosOperNa         &                    & \\
\bottomrule
\end{tabularx}
\end{table}

\subsection{RQ3 · Recuperación e instrumentación}
\label{sec:rq3}

Los embeddings son el mecanismo mayoritario (\clRQtresEmbeddings{}), seguidos del ranking léxico o
reordenamiento (\clRQtresRanking{}) y del recorrido de grafo (\clRQtresGrafo{}). La lectura directa
de archivos, sin índice, aparece en \clRQtresGrep{} estudios (\clRQtresGrepPct{}\,\%);
\clRQtresRetrievalNa{} no declaran mecanismo. Sobre la instrumentación —el núcleo de la
pregunta—, \clRQtresMide{} estudios (\clRQtresMidePct{}\,\%) miden de alguna forma qué se recuperó y
con qué utilidad, y \clRQtresNoMide{} no lo hacen. Conviene leer esa cifra con cuidado: «medir» aquí
abarca desde reportar la exactitud de recuperación hasta ablations por elemento, y la sección de
discusión muestra que la evidencia acumulada bajo esa etiqueta apunta en direcciones opuestas.

\subsection{RQ4 · Persistencia cross-sesión y multi-agente}
\label{sec:rq4}

El alcance dominante es el proyecto (\clRQcuatroProyecto{}), por delante de la sesión
(\clRQcuatroSesion{}), la tarea (\clRQcuatroTarea{}) y la organización
(\clRQcuatroOrganizacion{}): el objeto de estudio se desplazó efectivamente hacia la operación
continua. La memoria compartida entre agentes aparece en \clRQcuatroMultiagente{} estudios y la
verificación de fidelidad o procedencia en \clRQcuatroFidelidad{}. El traspaso explícito entre
sesiones o agentes está definido en \clRQcuatroHandoff{} (\clRQcuatroHandoffPct{}\,\%) y ausente en
\clRQcuatroHandoffNo{}; de los que lo definen, \clRQcuatroHandoffContenido{} traspasan contenido,
\clRQcuatroHandoffContrato{} un contrato y \clRQcuatroHandoffPunteros{} punteros. El estado del plan
de trabajo se trata como contenido de memoria en \clRQcuatroPlanEstado{} estudios, lo que confirma
como fenómeno real una dimensión que el protocolo había registrado sólo como hipótesis.

\subsection{RQ5 · Evaluación y costo}
\label{sec:rq5}

El corpus es empíricamente sólido en su dimensión principal: \clRQcincoBenchmark{} estudios evalúan
con un benchmark end-to-end y \clRQcincoAblation{} (\clRQcincoAblationPct{}\,\%) incluyen ablation
con y sin memoria, que es la evidencia causal que el mapeo buscaba. Sólo \clRQcincoCualitativa{}
aportan evidencia exclusivamente cualitativa. El costo se mide en \clRQcincoCosto{} estudios
(\clRQcincoCostoPct{}\,\%). El contraste está en el tipo de métrica: el éxito final se reporta en
\clRQcincoMetricaExito{} casos y el costo en tokens en \clRQcincoMetricaTokens{}, pero las métricas
de iteración aparecen sólo en \clRQcincoMetricaIteracion{} (\clRQcincoMetricaIteracionPct{}\,\%).

\begin{table}[!t]
\caption{RQ5 · Tipo de evidencia y métricas reportadas ($n=\clCarAnalizados$)}
\label{tab:rq5}
\centering
\begin{tabularx}{\columnwidth}{Xr@{\hskip 1.2em}Xr}
\toprule
\textbf{Tipo de evidencia} & \textbf{n} & \textbf{Métrica} & \textbf{n} \\
\midrule
Benchmark end-to-end & \clRQcincoBenchmark   & Éxito final    & \clRQcincoMetricaExito \\
Ablation con/sin     & \clRQcincoAblation    & Otra           & \clRQcincoMetricaOtra \\
Medición de costo    & \clRQcincoCosto       & Costo en tokens& \clRQcincoMetricaTokens \\
Sólo cualitativa     & \clRQcincoCualitativa & Latencia       & \clRQcincoMetricaLatencia \\
                     &                       & Iteración      & \clRQcincoMetricaIteracion \\
                     &                       & Ninguna        & \clRQcincoMetricaNinguna \\
\bottomrule
\end{tabularx}
\end{table}

\subsection{RQ6 · Riesgos y control}
\label{sec:rq6}

La fidelidad —procedencia, obsolescencia, consistencia entre almacenes— es una preocupación
extendida: \clRQseisFidelidadFaceta{} estudios la tocan. La seguridad \emph{de} la memoria, en
cambio, es marginal: \clRQseisSeguridad{} estudios (\clRQseisSeguridadPct{}\,\%). Entre los modos de
fallo nombrados explícitamente, la obsolescencia aparece en \clRQseisStaleness{} estudios, el
envenenamiento o inyección en \clRQseisPoisoning{} y la sobre-retención o el crecimiento sin control
en \clRQseisSobreRetencion{}. Un sistema que persiste conocimiento durante meses acumula
exactamente esos tres riesgos, y el corpus los cubre de forma desigual.

\subsection{Facetas transversales y lente de decisión}
\label{sec:res-facets}

La tabla~\ref{tab:facetas} ordena las diez facetas por prevalencia confirmada contra el texto
completo. La lectura importante es el extremo inferior, no el superior.

\begin{table}[!t]
\caption{Prevalencia confirmada de las facetas ($n=\clCarAnalizados$)}
\label{tab:facetas}
\centering
\begin{tabularx}{\columnwidth}{Xr}
\toprule
\textbf{Faceta} & \textbf{n} \\
\midrule
Recuperación medida            & \clMFretrievalMedido \\
Política de escritura          & \clMFwrite \\
Fidelidad y procedencia        & \clMFfidelidad \\
Costo y presupuesto de contexto& \clMFcosto \\
Consolidación y olvido         & \clMFconsolidacion \\
Conocimiento negativo          & \clMFnegativo \\
Memoria compartida multi-agente& \clMFmultiagente \\
Memoria procedimental          & \clMFprocedimental \\
Seguridad de la memoria        & \clMFseguridadFaceta \\
Documentos versionados         & \clMFdocsVersionado \\
\bottomrule
\end{tabularx}
\end{table}

La síntesis produjo \clMFentradas{} mapeos entre un estudio y una decisión de diseño concreta:
\clMFconfirma{} la confirman, \clMFextiende{} la extienden, \clMFneutral{} son neutrales y
\clMFrefuta{} la refutan, estos últimos repartidos en \clMFrefutaPapers{} estudios distintos. Por
decisión, los mapeos se concentran en la economía de contexto (\clMFanclaCosto{}), la fidelidad
(\clMFanclaFidelidad{}), la recuperación medida (\clMFanclaGMdos{}) y la consolidación
(\clMFanclaConsolidacion{}), seguidas del conocimiento negativo (\clMFanclaNegativo{}), las costuras
entre almacenes (\clMFanclaCosturas{}), el estado del plan (\clMFanclaPlanEstado{}), el disparador de
escritura (\clMFanclaGMuno{}), el índice como caché (\clMFanclaTierIndice{}) y los documentos como
fuente de verdad (\clMFanclaTierDocs{}). La refutación se concentra en la recuperación medida
(\clMFrefutaGMdos{}, \clMFrefutaGMdosPct{}\,\% del total) y, en segundo lugar, en el conocimiento
negativo (\clMFrefutaNegativo{}); ambas se analizan en la sección~\ref{sec:discussion}.
%%CLAIMS-STRICT-END

\subsection{RQ7 · Práctica industrial}
\label{sec:rq7}

% =====================================================================================
\section{Discusión}
\label{sec:discussion}

\subsection{El mapa de vacíos se invierte al medirlo}
\label{sec:disc-gaps}

%%CLAIMS-STRICT
Antes de cribar, el protocolo estimó la prevalencia de cada faceta con expresiones regulares sobre
títulos y resúmenes. Esa estimación situaba la política de escritura y el conocimiento negativo como
lo más escaso de la literatura, y de ahí se derivó una hipótesis de vacío. Medida contra el texto
completo de los \clCarAnalizados{} estudios, la estimación se invierte: la escritura está bien
estudiada (\clMFwrite{} estudios) y la escasez real está en otra parte. La faceta menos representada
del corpus es la memoria como \emph{documentos versionados} —archivos, control de versiones,
texto plano como sustrato— con \clMFdocsVersionado{} estudios, apenas
\clMFdocsVersionadoPct{}\,\% del corpus; la sigue la seguridad de la memoria con \clRQseisSeguridad{}
(\clRQseisSeguridadPct{}\,\%).
%%CLAIMS-STRICT-END

La lección metodológica es directa y aplica a cualquier mapeo que use pistas automáticas para
priorizar: una expresión regular sobre resúmenes mide \emph{vocabulario}, no prevalencia. La
calibración lo cuantificó antes de la pasada completa —de las coincidencias sugeridas para las dos
facetas más escasas, ninguna sobrevivió a la lectura— y por eso las pistas se usaron sólo para
ordenar la revisión, nunca para decidirla. Un mapeo que hubiera reportado los conteos de pista como
prevalencias habría publicado el vacío exactamente invertido.

\subsection{Una posición minoritaria bajo escrutinio}
\label{sec:disc-gmdos}

%%CLAIMS-STRICT
El sistema que motiva este mapeo tomó una decisión deliberada y documentada: recuperar leyendo
archivos directamente, sin índice semántico, tras medir que el índice no se pagaba en su escala. El
corpus muestra que esa es una posición minoritaria: \clRQtresGrep{} estudios
(\clRQtresGrepPct{}\,\%) recuperan por lectura directa, frente a \clRQtresEmbeddings{} que usan
embeddings. La pregunta interesante no es si la posición es popular, sino qué dice la evidencia que
la contradice.

De los \clMFrefuta{} mapeos que refutan alguna decisión del sistema, \clMFrefutaGMdos{}
(\clMFrefutaGMdosPct{}\,\% del total) apuntan a esta. Leídas una por una, no dicen lo mismo: se
reparten en dos direcciones opuestas y casi empatadas. \clGMdosMenos{} miden que \emph{añadir} una
capa de recuperación no ayuda o perjudica —sistemas de memoria comerciales sin ganancia significativa
sobre no tener memoria, un baseline de contexto largo sin memoria explícita obteniendo el mejor
rango, documentos completos en texto plano superando a recuperación vectorial y por grafo, un
puntuador determinista sin embeddings alcanzando precisión perfecta en el primer resultado sobre
documentos etiquetados—. Otras \clGMdosMas{} miden lo contrario: que la similitud \emph{pura} es
insuficiente, y que hace falta híbrido, estructura o recuperación agéntica multi-salto. Las
\clGMdosCalidad{} restantes no discuten recuperación sino que el número de titular del propio
trabajo contradice su tabla de resultados.
%%CLAIMS-STRICT-END

El empate no es ruido: es el hallazgo. La evidencia sobre recuperación medida está \emph{en disputa},
y cualquier revisión que cite sólo uno de los dos bloques está seleccionando. La reconciliación
razonable es que la decisión que se discute son en realidad dos, habitualmente confundidas en una:
(a)~si conviene tener una capa de recuperación, y (b)~con qué mecanismo. Sobre (a), la evidencia
favorece la abstención cuando el corpus cabe en la ventana y está curado —que es precisamente el
caso de un proyecto único con documentación versionada—. Sobre (b), la evidencia es netamente
desfavorable al \emph{léxico puro}: los trabajos que comparan mecanismos sobre la misma base miden
que el emparejamiento léxico es el más débil de los tres y que el vector denso domina salvo en
nichos de vocabulario cerrado, y que la similitud semántica sola falla en consultas por atributo
temporal.

La consecuencia práctica para un sistema como el estudiado no es «agregar embeddings», sino
reconocer que su decisión es defendible por el argumento (a) y no por el (b), y que su condición de
validez es explícita y comprobable: que el corpus siga cabiendo en la ventana y que las consultas
sigan siendo de recuperación estructurada y no de similitud abierta. El corpus también entrega el
instrumento con el que esa condición se vigila —la medición de utilidad por lectura— y muestra que
\clRQtresMide{} estudios (\clRQtresMidePct{}\,\%) ya la instrumentan de alguna forma.

\subsection{Qué debería cambiar en el diseño de agentes de larga vida}
\label{sec:disc-implications}

%%CLAIMS-STRICT
Tres asimetrías del corpus tienen consecuencia de diseño inmediata.

\textbf{El disparo de escritura ya no es discrecional.} Entre los estudios que declaran qué dispara
la escritura, \clRQdosDeterministaPct{}\,\% usan un mecanismo no discrecional: evento determinista
(\clRQdosTriggerDeterminista{}), checkpoint (\clRQdosTriggerCheckpoint{}) o presupuesto
(\clRQdosTriggerPresupuesto{}), frente a \clRQdosTriggerDiscrecional{} que dejan la decisión al
criterio del modelo. La práctica de «el agente escribe cuando le parece relevante» es minoritaria en
la literatura, y quien la conserve debería justificarla.

\textbf{El traspaso entre sesiones es norma, pero viaja mal.} \clRQcuatroHandoff{} estudios
(\clRQcuatroHandoffPct{}\,\%) definen un traspaso explícito entre sesiones o agentes, de modo que la
discontinuidad ya se reconoce como el problema central. Ahora bien, casi todos traspasan
\emph{contenido} (\clRQcuatroHandoffContenido{}); traspasar \emph{punteros}
(\clRQcuatroHandoffPunteros{}) o un \emph{contrato} de lo que debe cumplirse
(\clRQcuatroHandoffContrato{}) sigue siendo raro. Es un espacio de diseño casi vacío, y es
justamente el que evita re-copiar contexto que envejece.

\textbf{El campo mide si funciona, no si deja de repetir errores.} La evidencia causal es sólida:
\clRQcincoAblation{} estudios (\clRQcincoAblationPct{}\,\%) incluyen ablation con y sin memoria. Pero
las métricas de iteración —pasos, reintentos, errores repetidos, hipótesis recicladas— aparecen en
\clRQcincoMetricaIteracion{} (\clRQcincoMetricaIteracionPct{}\,\%). El mapeo previo del equipo ya
había señalado esta brecha con veintisiete estudios; se confirma sobre \clCarAnalizados{}. Un agente
de larga vida falla repitiéndose, y ésa es exactamente la dimensión que casi nadie mide.

Vale la pena registrar además dónde el corpus está desierto: sólo \clRQunoRepParametros{} estudios
escriben la memoria en los parámetros del modelo, y \clRQunoSinMecanismo{} de los analizados no
aportan mecanismo propio sino un benchmark, un dataset o una medición. Esto último es un signo de
madurez incipiente: ya hay casi tantos protocolos de evaluación como sistemas que evaluar.
%%CLAIMS-STRICT-END

% =====================================================================================
\section{Amenazas a la validez}
\label{sec:threats}

\subsection{Validez de constructo}

«Memoria» es un término sobrecargado y el estudio depende de dónde se trace su borde. La rúbrica
exige que el mecanismo sea \emph{persistente} —entre tareas o entre sesiones— o \emph{compartido}
entre agentes; en consecuencia, el manejo intra-tarea de la ventana de contexto de un solo agente
queda fuera aunque sus autores lo llamen memoria. Esa frontera es defendible dado el ancla del
estudio, pero es una decisión, no un hecho: un lector interesado en compresión de contexto encontrará
que trabajos relevantes para él fueron excluidos por diseño. El vocabulario de facetas es cerrado y
se congeló antes de la pasada completa; una dimensión que ninguna faceta nombra es invisible para
este mapa.

\subsection{Fiabilidad de la clasificación}

%%CLAIMS-STRICT
Ésta es la amenaza principal y conviene enunciarla sin adornos. La clasificación la realizaron
\emph{modelos de lenguaje} bajo una rúbrica congelada, con compuertas deterministas que verifican la
aritmética, los vocabularios y la consistencia entre fases, pero que no pueden verificar el juicio.
La segunda pasada descrita en la sección~\ref{sec:reliability} es un segundo lector \emph{modelo},
ciego a la decisión del primero: mide consistencia entre lecturas independientes, no acuerdo entre
dos personas expertas. Dos lectores que comparten arquitectura pueden compartir sesgos, y ese
componente de error es invisible desde dentro del propio experimento; el coeficiente reportado debe
leerse como una cota de consistencia interna del procedimiento, no como validación externa. La
adjudicación humana se aplicó sobre los conflictos, no sobre el corpus completo.
%%CLAIMS-STRICT-END

\subsection{Cobertura y sesgo de accesibilidad}

%%CLAIMS-STRICT
La búsqueda usó tres agregadores con exportación programática y quedó fuera la literatura indexada
únicamente en bases con licencia. Un control de sanidad con obras validadas por la comunidad y con
el corpus del mapeo previo mostró que las ausencias corresponden a un hueco de \emph{fuente} —obras
no indexadas en ninguno de los tres agregadores— y no de vocabulario, y esas obras se recuperaron por
snowballing.

Más importante, y menos habitual de reportar: la pérdida de texto completo \textbf{no es aleatoria}.
De los \clSelIncluidosTotal{} incluidos se analizaron \clSelConTexto{} (\clSelCobertura{}\,\%); los
\clSelSinTexto{} restantes quedaron sin PDF por protección anti-robots, ausencia de acceso abierto o
páginas intermedias. De ellos, exactamente \clAmeSinTextoPreprint{} son preprints: el material
inaccesible proviene íntegramente de venues con DOI de editorial, mientras que
\clAmeAnalizadosPreprint{}\,\% de lo analizado son preprints. El corpus analizado está por tanto
sesgado hacia lo libremente accesible, y ese sesgo apunta en la misma dirección que la señal de
calidad de la próxima subsección.
%%CLAIMS-STRICT-END

\subsection{Calidad del corpus}

%%CLAIMS-STRICT
El corpus es muy reciente —\clCarRecientesPct{}\,\% de los últimos dos años— y por tanto contiene una
proporción alta de preprints sin revisión por pares. La digestión marcó \clCalReconsiderarDig{}
estudios (\clCalReconsiderarDigPct{}\,\%) con señales de calidad dudosa: venues autopublicados,
evidencia sólo cualitativa presentada como medición, contradicciones entre el resumen y las tablas
del propio trabajo y, en algunos casos, texto que parece generado por un asistente en lugar de
resultados. La extracción, que leía campos en vez del texto completo, había marcado
\clCalReconsiderarExt{}: leer el documento entero detecta más del doble de problemas. Ninguno de esos
registros se descartó automáticamente —se marcaron y quedaron sujetos a revisión—, porque una señal
de sospecha no es un criterio de exclusión. El lector debe saber que un corpus de esta cosecha
incluye trabajo de calidad muy desigual, y ese hecho es en sí un resultado sobre el estado del área.
%%CLAIMS-STRICT-END

\subsection{Conflicto de interés analítico}

La lente que clasifica cada estudio como confirmación o refutación de una decisión de diseño la
aplica el mismo equipo que tomó esas decisiones. Es un riesgo evidente de complacencia. Se mitigó de
tres formas comprobables: el vocabulario de posturas es cerrado y la refutación se declaró de
antemano como la categoría de mayor valor; la clasificación por dirección de las refutaciones sobre
la decisión más disputada vive en un archivo aparte, con la evidencia citada de cada asignación; y
una compuerta exige que esa clasificación cubra \emph{exactamente} el conjunto de refutaciones del
corpus, de modo que descartar en silencio una fila incómoda hace fallar la compilación. Aun así, la
mitigación es procedimental: no sustituye a una réplica independiente.

\subsection{Reproducibilidad}

Protocolo, hojas de decisión, formularios, veredictos por lote y el registro de afirmaciones están
versionados; cada cifra de este artículo se re-deriva de esos artefactos en cada compilación, de modo
que un tercero puede recomputarlas o detectar que dejaron de seguirse del corpus. Los PDF no se
distribuyen por razones de licencia: se publica su índice con identificador y hash criptográfico, que
permite verificar que el documento analizado es el mismo sin redistribuirlo.

% =====================================================================================
\section{Conclusiones}
\label{sec:conclusions}

% =====================================================================================
\bibliographystyle{IEEEtran}
\bibliography{sms-v2}

\end{document}
